% !TEX program = lualatex
% Build: latexmk -lualatex syllabus.tex   (or: lualatex syllabus.tex)
\documentclass[11pt,a4paper]{article}

\usepackage[margin=25mm]{geometry}
\usepackage{luatexja}
\usepackage[haranoaji]{luatexja-preset}
\usepackage{enumitem}
\usepackage{booktabs}
\usepackage{tabularx}
\usepackage{xcolor}
\usepackage{titlesec}
\usepackage[hidelinks]{hyperref}
\usepackage{fontawesome5}

\setlist{nosep,leftmargin=*}
\setlength{\parindent}{0pt}
\setlength{\parskip}{0.6em}
\titleformat{\section}{\large\bfseries}{}{0pt}{}
\titlespacing*{\section}{0pt}{1.4em}{0.4em}
\renewcommand{\arraystretch}{1.25}

% Items to be confirmed before distribution.
\newcommand{\tbd}[1]{\textcolor{red}{[#1]}}
\newcolumntype{Y}{>{\raggedright\arraybackslash}X}
\newcommand{\extlink}[2]{\href{#1}{\textcolor{blue}{#2\,\textsuperscript{\faLink}}}}

\begin{document}

\begin{center}
  {\LARGE\bfseries 英語原書講読入門}\\[0.3em]
  {\Large\bfseries Introduction to Academic Reading in English}\\[0.6em]
  {\large Syllabus \quad Fall 2026}
\end{center}

\section{Course Information}
\begin{tabularx}{\textwidth}{@{}lY@{}}
  \toprule
  Instructor & Zeyu Lyu \\
  Meeting time & Friday  13:00--14:30  \\
  Location & C403 \\
  Contact & lyu.zeyu.e8@tohoku.ac.jp \\
  Main text & Matthew J.\ Salganik, \emph{Bit by Bit: Social Research in the Digital Age}
              (Princeton University Press, 2018). 
              \extlink{https://www.bitbybitbook.com/en/1st-ed/preface/}{Link} \\
  Japanese edition & マシュー・J.\ サルガニック(著)、瀧川裕貴・常松淳・阪本拓人・大林真也(訳)
              『ビット・バイ・ビット――デジタル社会調査入門』有斐閣、2019年
              (\extlink{https://www.yuhikaku.co.jp/book/b10189273.html}{Link}) \\
  \bottomrule
\end{tabularx}


\section{Course Objectives}
\begin{enumerate}
  \item Learn how to read and understand academic texts in English.
  \item Learn how to find and use relevant background information to better understand academic texts in English.
  \item Learn how to summarize and present academic texts in English.
  \item Develop a basic understanding of computational social science and its
        research methods.
\end{enumerate}

\section{Assessment}
\begin{itemize}
  \item Attendance and class participation: 70\%
  \item Group presentation: 30\%
\end{itemize}

\newpage
\section{Schedule}

\begin{tabularx}{\textwidth}{@{}rlYl@{}}
  \toprule
  \# & Date & Topic & Format \\
  \midrule
   1 & Oct 2  & Introduction & \\
   2 & Oct 9  & Introduction (Chapter 1) Observing behavior I (Chapter 2) & \\
   3 & Oct 30 & Observing behavior II (Chapter 2) & \\
   4 & Nov 6  & Asking questions I (Chapter 3) & \\
   5 & Nov 13 & Asking questions II (Chapter 3) & \\
   6 & Nov 20 & Running experiments I (Chapter 4) & \\
   7 & Nov 27 & Running experiments II (Chapter 4) & \\
   8 & Dec 4  & Creating mass collaboration I (Chapter 5) & \\
   9 & Dec 11 & Creating mass collaboration II (Chapter 5) & \\
  10 & Dec 18 & Ethics (Chapter 6) and The future (Chapter 7) & \\
  11 & Dec 25 & Presentations & Online \\
  12 & Jan 8  & Presentations & Online \\
  13 & Jan 15 & Presentations & Online \\
  14 & Jan 22 & Presentations & \\
  15 & Jan 29 & Presentations & \\
  16 & Feb 3 (Wed) & Reserve session & \\
  \bottomrule
\end{tabularx}

\small\emph{Note: Class dates and the course schedule are subject to change.}
\normalsize

\section{Course Content}

\subsection{Reading the English Texts}

\begin{itemize}
  \item The main text is written in English. Students are expected to read the
        assigned chapters in English before each class. Students are encouraged
        to use dictionaries and other resources to help them understand the text.
  \item In each class, the instructor will explain the main points of the
        assigned chapter and discuss any questions or difficulties that students
        may have. Students are encouraged to ask questions and participate in
        the discussion.
\end{itemize}

\subsection{Presentations}

\begin{itemize}
  \item Groups will normally consist of two or three students; the final number
        of groups will be determined after enrollment is finalized.
  \item Each group will choose one paper from a reading list provided by the
        instructor.
  \item Each group will have 30--45 minutes for its presentation and the
        discussion that follows. Two groups will present in each session.
  \item The presentation order will be announced on December 4. Please form
        your group by that date.
\end{itemize}

\end{document}
